% !PRINTED 489
\weberpage{489}

Daraus erhält man

\[
\begin{aligned}
u^3v^3 &= 4\,(3+\sqrt{13}), \\
u^6+v^6 &= 4\,(17+5\sqrt{13}), \\
u^3-v^3 &= \sqrt{2}\,(3+\sqrt{13}),
\end{aligned}
\]
so daß, wenn $u^3 = \sqrt{8}\,x$ gesetzt wird, für $x$ die quadratische Gleichung folgt:

\weq{17}{x^2 - \frac{3+\sqrt{13}}{2}\,(x+1) = 0, \quad f(\sqrt{-39})^3 = \sqrt{8}\,x.}

Wir wollen unsere Resultate jetzt noch anwenden auf zwei Probleme, die in die allgemeine Transformationstheorie des siebenten Abschnittes gehören.

\begin{center}
\textbf{\S\,132. \gesperrt{Die Schlaeflische Modulargleichung für den 23sten Transformationsgrad.}}
\end{center}

Jede Klassengleichung tritt als Divisor in einer großen Zahl von Transformationsgleichungen auf und kann daher, wenn sie auf andere Weise bekannt ist, zur Berechnung von Transformationsgleichungen benutzt werden. Wir nehmen als Beispiel den 23sten Transformationsgrad.

Nach \S\,73 besteht, wenn

\weq{1}{u = f(\omega), \quad v = f(23\,\omega), \quad f\left(\frac{c+\omega}{23}\right), \quad c \equiv 0 \ (\mathrm{mod}\ 48)}

gesetzt ist, eine Gleichung zwischen

\weq{2}{\begin{aligned}
A &= \left(\frac{u}{v}\right)^{12} + \left(\frac{v}{u}\right)^{12} \\
B &= uv + \frac{2}{uv},
\end{aligned}}

und es ist schon in \S\,74, (14) gezeigt, daß diese Gleichung die Form hat:

\weq{3}{A = B^{11} + m_1 B^{10} + \cdots + m_{10} B + m_{11},}

worin die Koeffizienten $m_1, m_2 \ldots, m_{11}$ rationale Zahlen sind. Statt nun diese rationalen Zahlenkoeffizienten wie dort aus den Entwickelungen von $u$ und $v$ nach Potenzen von $q$ zu berechnen, suchen wir die Wurzeln der Gleichung (3) für $u=v$ aus der komplexen Multiplikation zu bestimmen.
